\subsubsection{Trace reductions for rectangular quartics}
\label{app:rectangular-trace-details}

We complete the trace calculations for
\cref{prop:rectangular-minimal-selection}, using
$B=X^{\mathsf T}X$, $P_j=\tr(B^j)$, and $R=P_1$.
In \eqref{eq:rectangular-trace-differentiation}, summing coordinate
derivatives of $XB$ gives $(p+3+1)P_1$. For $XB^2$, the derivative
of the first factor contributes $pP_2$; the two derivatives of the
first $B$ contribute $P_2$ and $P_1^2$, and those of the second
contribute $P_2$ and $3P_2$. Thus
\[
 \Delta P_2=4(p+4)P_1,\qquad
 \Delta P_3=6\bigl((p+5)P_2+P_1^2\bigr).
\]
These identities hold without any regular-eigenvalue assumption.

For the three eigenvalues of $B$, put
\[
 e_2=(R^2-P_2)/2,\qquad
 e_3=(P_3-RP_2+Re_2)/3.
\]
Their characteristic equation gives
$P_j=RP_{j-1}-e_2P_{j-2}+e_3P_{j-3}$ for $j\ge4$.
Together with \eqref{eq:rectangular-S-traces}, this yields
\begin{align*}
 P_2&=(F+R^2)/2,\\
 P_3&=(S+32RF+16R^3)/64,\\
 P_4&=RS/48+5R^2F/12+R^4/8+F^2/8,\\
 P_5&=5R^2S/256+5R^3F/16+R^5/16
                 +5RF^2/24+5FS/768.
\end{align*}
Differentiation of $S$ gives
\[
 \nabla S=X\bigl(384B^2-256RB+(32R^2-64F)\Id\bigr).
\]
Pairing this with $\nabla F=4X(2B-R\Id)$ and with itself, and then
substituting the displayed power traces, proves the last two entries
of \eqref{eq:rectangular-quartic-gamma}.
